\section{Deep values: most people share some degree of universalism}



% 2. Deep values: most people share some degree of universalism

% =======================================
% - general questions (ISSP, UNDP, cf. A.1.1 de mon papier NHB)
% =======================================

ISSP surveys indicate that concern for global problems is displayed by a majority. Across 36 countries, 82\% of respondents agree that environmental problems require international agreements that countries must follow, while only 4\% disagree. In a more recent survey covering 29 countries, 78\% respondents consider existing economic differences between rich and poor countries too large, while 5\% who disagree \citep{issp_environment_2010,issp_inequality_2019}. 

% =======================================
% - wrong perceptions (Fehr et al 2022)
% =======================================
Despite such concerns, there is a gap between attitudes toward global inequalities and actual policies. A first attempt to explain it comes from misperception. Indeed, Germans underestimate their global income rank by an average of 15 percentage points. Correcting perceived national rank changes, both national and global, redistributive preferences among left-leaning respondents, whereas correcting perceived global rank has no detectable effect on support for global redistribution. In addition, 90\% of respondents expressed a preference for some degree of global redistribution, although intensity vary across individuals.
By contrast, aid support in US is driven by information on global distribution, because people underestimate their rank by 27 percentage points in their global income rank and overestimate global median income by a factor 10, they tend to increase their support to foreign aid by almost 10 percentage points and donation to international charities by 55\% \citep{nair_misperceptions_2018}.

% =======================================
% - universalism (Enke, etc.), cf. A.1.5 
% =======================================

Literature has long studied the concept of global identity (\citet{rey_2018} for a review). Universalism has several definitions in the literature, although some are extremely close. \citet{enke_moral_2023}, \citet{enke_universalism_2024}, and \citet{cappelen_universalism_2025} define it as the extent to which altruism remains invariant to recipients' identity, group membership, or social and geographic distance. \citet{boyer-kassem_how_2026} retain this conceptual definition but argue that its measurement requires recipients to be economically comparable and to belong to clearly identified groups. A related concept is the moral circle \citep{waytz_ideological_2019,jaeger_toward_2026}.

\citet{bayram_2015} indicates that ``78\% of the participants in 57 countries see themselves as citizens of the world'', referring to the 2005-2008 wave of the World Values Survey, while the 2017-2022 wave % AF: link to https://www.worldvaluessurvey.org/WVSOnline.jsp
% exposes a closer feeling to their town, region or country compared to the world. 
show that 84\% feel close (or very close) to their country but only 47\% feel close to the world. 
\citet{nation_2024} argues that across 21 countries, there is a large variation since 31\% to 88\% of respondents see themselves ``more a world citizen than a citizen of $\left[ \text{their} \right]$ country''. The cross-country average is about 46\% while this share exceeds 50\% in seven countries, with Indonesia and India exhibiting the highest level, and the Netherlands and China displaying the lowest\footnote{Calculated excluding non-committal responses (``don't know'', ``prefer not to say'', and ``neither agree nor disagree''). The underlying country-level data were obtained courtesy of the authors.}. % While comparable shares agree on earmarking tax revenue for global issues. AF: this sentence is unclear; if you refer to new results, give them more precisely (with figures)

Universalism is reflected not only in stated preferences but also in observed patterns of giving and redistribution. \citet{enke_universalism_2024} measure universalism at the U.S. district level using donation data, and find that a district's universalism, defined as the extent to which charitable giving is invariant to geographic and social distance, predicts electoral outcomes better than its income or education level. Authors % AF: c'est le même papier qui a les deux méthodos ?
measure universalism at the individual level asking American respondents to split \$100 between a random stranger and a random compatriot %individual 
with the same income.% but closer
\footnote{These results nevertheless faces measurement concerns as pointed out by  \citet{boyer-kassem_how_2026}. Authors argue that equal nominal incomes do not neutralize differences in purchasing power, wealth, or relative social status, and that a randomly selected ``stranger'' may itself belong to the relevant in-group, potentially confounding moral distance with economic considerations.} % AF: la footnote n'est pas clair (surtout la fin)
They use different angles of universalism and define foreign universalism as the tendency to give to a foreigner preferably to a fellow citizen. Authors note that most % AF: c'est quoi "most"? il faut des chiffres
people are biased in favour of their own country, which could be explained in part by concerns about inequality, since the split involves two people on the same nominal income, with the foreigner almost surely living in a poorer country than the American and therefore benefits from a higher social status. 

\citet{prather_2013} controls for inequality and still finds a nationalist bias in the United States, where 84\% of Americans agree that ``taking care of problems at home is more important than giving aid to foreign countries'' \citep{PIPA2001}. Support for a U.S. program to fight hunger is 10 percentage points lower when the hypothetical program operates abroad rather than domestically, and the gap reaches 20 percentage points when the choice concerns cutting an existing program.

Using a similar approach, \citet{enke_moral_2023} show that moral universalism % AF: comment ils le définissent ?
explains political views better than standard political economy variables such as income, wealth or perception about government efficiency across Australia, Germany, France, the US, and Sweden. For instance, if universalism levels were the same for all, individuals' policy position would not be as much polarized as they are. Similarly, using surveys covering 60 countries and 64,000 respondents, \citet{cappelen_universalism_2025} document substantial individual variation in universalism: 26\% of respondents divide resources equally in every allocation task, whereas 17\% allocate at most 20\% to strangers across the different tasks. % AF: les tâches consistent en quoi ? Quelle est la proportion d'universaliste ? Comment ça diffère par pays ?

Universalism can be surprising because it may lead people to sacrifice their own well-being for the sake of distant groups of people. This is particularly true of climate policies % AF: pourquoi climate alors que l'exemple qui suit ne parle pas de climat ?
that are implemented at local level but have global effects. In a lab experiment with U.S. students, \citet{cherry_2017} find that 61.1\% of participants support egalitarian policies even when these reduce their own material payoff. Complementary international experimental evidence from a lab experiment matching participants from Norway and Germany with participants from Uganda and Tanzania in a dictator game shows that participants acted as moral cosmopolitans and did not assign importance to nationality in their distributive choices \citep{cappelen_needs_2013}. Likewise, \citet{leiserowitz_2006} reports that a majority of Americans express greater concern about the consequences of climate change for people worldwide (50\%) or for non-human nature (18\%) than for themselves and their families (12\%) or for the United States (9\%). Consistent with these findings, a 2017 Focus 2030 survey shows that 40\% of French respondents agree that reducing poverty in developing countries should be a priority for the European Union, compared with only 19\% who disagree.

% =======================================
%- moral circle, group defended: Waytz et al 19, Jaeger & Wilks 23 (give descriptive stats and not only correlation if possible), Fabre et al 25, Fabre 26
% =======================================

A complementary way to describe this heterogeneity is through the concept of the ``moral circle'' defined as the set of people or other entities considered worthy of moral consideration. Across seven study samples with U.S. respondents, \citet{waytz_ideological_2019} show that left-leaning individuals consistently report broader moral circles than conservative ones. % AF: il ne suffit pas de dire ça, il faut donner les stats descriptives pour l'ensemble de la population, puis pour gauche vs. droite
In two studies in the Netherlands, US, UK and Australia, \citet{jaeger_relative_2023} find that differences between the groups being evaluated explain 40\% of the variation in moral concern, differences between respondents explain 29\%, and their interaction explains 27\%. % AF: expliquer mieux ce que c'est que "differences between groups", la phrase n'est pas claire

We also find this structure when respondents are asked which group they advocate for when they vote. Among 8,000 respondents in five Western countries, \citet{fabre_majority_2025} argue that ``When we ask respondents which group they defend when they vote, 20\% choose ‘sentient beings (humans and animals)’, 22\% choose ‘humans’, 33\% select their ‘fellow citizens’ (or ‘Europeans’), 15\% choose ‘My family and myself’ and the remaining 9\% choose another group (mainly ‘My State or region’ or ‘People sharing my culture or religion’). Notably, a majority of left-wing voters choose humans or sentient beings.''. In a broader survey of 12,000 respondents across eleven 
high-income 
countries, \citet{fabre_public_2026} uses a new question to measure universalism asking respondents the group they advocate when voting and ``45\% choose a universalist  response (either “Humans” or “Sentient beings (humans and animals)”), while 32\% opt for their fellow citizens''. The universalist share reaches 50\% in Europe, 57\% in Saudi Arabia, and 59\% among left-wing respondents \citep{fabre_public_2026}. Universalistic concern is widespread, despite median voters favoring their fellow citizens in the U.S. and Japan.

% =======================================
%- why HICs should help LICs: duty (Fabre 26)
% =======================================


When respondents are asked why ``high-income countries should support low-income countries'', they first invoke a moral duty, on average at 54\% at the exception of France (43\%) and Japan (45\%), while, on average, 38\% of respondents suggest the ``Long-term interest of HICs to help LICs''. Finally, only 16\% of respondents advocate non of these arguments, which puts an end to self-interest or guilt bias \citep{fabre_public_2026}.

